\mathplain
\section{Conclusion}\label{sec:conclusion}

The textbook proof of type safety inverts subtyping, and in DOT subtyping cannot be inverted in every context, because a declaration in the context is an assumption that may be false. The proof in this report does not invert subtyping. It makes every use of subtyping an explicit proof term, and over a typed store, where every assumption is discharged by a stored literal, it normalizes those proof terms instead. Three design decisions make the normalization a structural induction: types refer to terms only through the block names of binders, stores hold literals at their precise types, and object coercions are templates that read named propositions of their source and adjust them by closed evidence. Preservation and progress then have the textbook shape, and the safety of DOT is a corollary that needs no metatheory of DOT beyond erasure.

\paragraph{Does the proof keep its shape?}
We expected it to. A new sort of member adds constructors to types and propositions, evidence rules for them, and cases to the translation. The normalizer gains cases and the canonical forms theorem gains cases, while preservation and progress only ask that theorem for a shape. Four extensions have since been built, each as a separate copy of the development: paths with singleton types, capture checking in two versions, and capture classifiers. Each copy changed the target, the normalizer, and the canonical forms theorem, and preservation and progress grew as well. Each carries the same theorems: canonical forms of closed evidence, preservation, progress, the erasure simulation, and the safety of its source by translation. None needed tight typing, a pushback measure, or a semantic model. The copies are separate developments, so this is evidence for the expectation and not a theorem about it.

\paragraph{Evidence from a compiler.}
A small surface language with a derivation-producing type checker exists as well, as Section~\ref{sec:mechanization} describes. The use we expect to matter most is still open. A compiler for a language in the DOT family could emit evidence for the subsumption steps it takes, in the way GHC emits FC coercions, and have them validated independently of its own subtyping algorithm. Validation is decidable, as Section~\ref{sec:related} notes, while subtyping is not, so the checker never has to search. That helps most for the parts of a compiler that have no reference semantics beyond a paper, and capture checking is one of them.
